\documentclass[12pt]{article}
\usepackage{amsmath,amsfonts,fullpage,amsthm,graphicx}

\newtheorem*{theorem}{Theorem}
\newtheorem*{fact}{Fact}
\newtheorem*{goal}{Goal}
\newtheorem*{idea}{Idea}
\newtheorem*{defn}{Definition}
\newtheorem*{ex}{Example}

%Style section
%\setlength{\textheight}{10in}
%\setlength{\textwidth}{7in}
%\hoffset=-0.75in
 \pagestyle{plain}
% \setlength{\topmargin}{0in}
% \setlength{\headsep}{0in}
% \setlength{\oddsidemargin}{0in}
% \setlength{\evensidemargin}{0in}

\begin{document}


\pagestyle{empty} %use this to get rid of page numbers

\noindent
{Math 166 \hskip 1.5 truein  Names:\ \hrulefill}

\noindent
%\vskip 0.3truein
%\bigskip
%\noindent
{Volumes by rotation again!}

%\noindent  \hskip 2.8 truein  Section Number:\ \hrulefill
\medskip

\begin{goal} \rm
Use integration to find the volume of various \emph{solids of revolution}.
\end{goal}
%\smallskip


In the questions 1 and 2, first graph the function(s). Determine whether to integrate with respect to $x$ or $y$ for the disk/washer method. State whether the cross-section is a disk or a washer. Then set up the integral(s). 
Then determine whether to integrate with respect to $x$ or $y$ for the shell method. Then set up the integral(s). 

(Hint: You will sometimes need to use two integrals.)

\begin{enumerate}
	
\item Find the volume of the solid obtained by revolving the region enclosed by $y = 32- 2x$, $y = 2x + 4$, and $x = 0$ about the $x$-axis.

\vspace{3in}

\item Now find the volume of the solid obtained by revolving that same region about the $y$-axis. 

\vspace{4in}

\item Use the disk or washer method to find the volume of the solid obtained by revolving the region enclosed by $y = x^2$, $y = 12 - x$, and $x = -1$ about the line $y = 15$. (Hint: Use your graph to figure out what your radius needs to be since we are revolving about the line $y = 15$ instead of the $x$-axis.)

\vspace{4in}

\item Use the shell method to find the volume of the solid obtained by revolving the region enclosed by $y = x^2$, $y = 12 - x$, and $x = -1$ about the line $x = -2$.

\newpage
%In the questions 3 and 4, first graph the function(s). Determine whether to integrate with respect to $x$ or $y$. Then set up the integral(s). Integrate if you have time. 
%\item Sketch the hyperbola $y^2-x^2=1$ by plugging in $-2,1,0,1,2$ for $x$ and solving for $y$. (Hint: there will be two $y$ values for each $x$ value.) Find the volume of the solid obtained by revolving the region between the branches of the hyperbola for $-2\leq x\leq 2$ about the $x$-axis. Note the same solid would result if you revolved the region bounded by the top branch of the hyperbola, the lines $x=2$ and $x=-2$, and the $x$-axis, about the $x$-axis. (This solid is called a \emph{hyperboloid}.)
\item An underground hemispherical dome with base a circle of radius 20 m is filled with mud of density 1906 kg/$m^3$. (Note the dome  is oriented so the flat part is farthest below ground.) Find the work done against gravity to pump enough mud out of the dome so that the level of the mud in the dome is 2 m after pumping. The dome is buried 1 meter below ground and the mud travels up the pipe and is dispensed on top of the ground. Draw a relevant picture in the coordinate axes and explain each step of your work.

%\begin{comment}
\begin{enumerate}
\item First draw a picture of the dome in the coordinate axes. 
%with the $y$-axis as the vertical axis and the origin at the bottom of the tank. 
Draw a horizontal layer at height $y$.
\vfill\vfill
\item Next, find the area $A(y)$ of the layer at height $y$ as a function of $y$.
\vfill
\item Write down an expression for the volume of this layer if it has thickness $\Delta y$.
\vfill
\item What is the mass of this layer?
\vfill
\pagebreak
\item What is the force due to gravity on this layer? 
\vfill
\item What is the vertical distance that this layer is lifted?
\vfill
\item Use your answers to (e) and (f) to find the work needed to pump the mud in this layer out of the spout.
\vfill
\item Integrate this expression to solve the problem.
\vfill
\end{enumerate}
%\end{comment}

\end{enumerate}

\end{document}
